\chapter*{Résumé}

\areprendre{
Ce projet de thèse porte sur l'étude de méthodes de discrétisation pour la décomposition de Helmholtz-Hodge dans le cas d'une surface fermée de $\R^3$ sans bord et leur application potentielle en électromagnétisme pour l'EFIE (Electric Field Integration Equation). On se propose d'adapter des méthodes volumes finis à grilles décalées de type Discrete Duality Finite Volume (DDFV) ayant déjà été employées dans le cas de domaines à bord dans $\R^2$. Cette adaptation nécessite de résoudre des difficultés posées par la géométrie du problème et qui affectent tant l'implémentation numérique des méthodes que l'étude théorique des schémas. Les propriétés intrinsèques des opérateurs DDFV, que ce soit en terme d'identités vérifiées au niveau discret ou bien de dimension des espaces présentent un intérêt notamment dans les méthodes de préconditionnement utilisées en EFIE.
}

\bigskip

\begin{envide}[Mots-clés]
Méthodes de volumes finis, Analyse Numérique, Électromagnétisme, Décomposition de Hodge-de Rham
\end{envide}

% Titre ayant l'apparence d'un chapitre sans saut de page
\begingroup
\makeatletter
\chaptermark{Abstract} % si vous voulez mettre à jour les en-têtes
\@makeschapterhead{Abstract}
\makeatother
\endgroup

\areprendre{
This thesis project focuses on the study of discretization methods for the Helmholtz-Hodge decomposition in the case of a closed surface in $\R^3$ without a boundary, and their potential application in electromagnetics to the Electric Field Integration Equation (EFIE). We propose to adapt Discrete Duality Finite Volume (DDFV) methods with staggered grids, which have already been used for bounded domains in $\R^2$. This adaptation requires resolving difficulties posed by the problem’s geometry, which affect both the numerical implementation of the methods and the theoretical analysis of the schemes. The intrinsic properties of DDFV operators, whether in terms of identities verified at the discrete level or the dimension of the spaces, are of particular interest in the preconditioning methods used in EFIE.
}

\bigskip

\begin{envide}[Keywords]
Finite Volumes methods, Numerical Analysis, Electromagnetism, Hodge-de Rham decomposition
\end{envide}